\chapter*{Preface}
\addcontentsline{toc}{chapter}{Preface}

This textbook is written for undergraduate students taking a first university course in cybersecurity, and for practitioners who wish to consolidate their understanding of the field on a sound conceptual basis. It assumes basic familiarity with computers and networks, but no prior study of security.

The book is organised into thirteen chapters grouped in four parts. The sequence is deliberate. It begins with concepts and principles, moves to the operating system as the first concrete security boundary, and then studies the adversary: who attacks, with which tools and through which channels. Only after the reader understands what must be defended does the book introduce the protective mechanisms—cryptography, identity management, secure software engineering and testing. The final part turns to operations: detecting and responding to incidents, analysing malicious code and evidence, and governing security as an organisational function.

Each chapter opens with a short introduction and a set of learning outcomes, develops its subject in numbered sections, and closes with a glossary of key terms, a summary and review questions. Tables condense comparisons that are easier to grasp side by side, while boxed examples connect abstract ideas to realistic situations. Cross-references between chapters are frequent, because security is a connected discipline: a weakness in one layer is usually exploited through another.

The book is published in parallel English and Greek versions. The Greek text is not a literal translation; it is an editorial rendering that follows the conventions of Greek academic writing and uses consistent, established terminology. Where no widely accepted Greek term exists, the English term is given in parentheses on first use.
